\begin{proposition}[E.37 --- Proved: an integral action detector and a zero-abelian obstruction at nineteen]
\label{E-integral-detector}
Retain the cubic lattice $A=\operatorname{Lie}_3(V)$,
the genuine normalized cocycle $\eta$ of E.19,
and the unit tangent at $O$ used in E.36.
There is an explicit surjective integral
$G_{\mathbb Q}$-module map
\[
 \Psi:A(-1)\longrightarrow\mathbb Z_3(1)
 \quad\text{with}\quad
 \Psi\eta(g)=6\kappa_3(3)(g)
 \quad(g\in G_{\mathbb Q}).
 \tag{E.37.1}
\]
The Kummer cocycle is normalized using the
positive real compatible roots of $3$.
In particular the class $[\eta]$ is nonzero
already after the displayed coefficient
pushout and reduction modulo nine.

The actual zero-abelian class-two section
\[
 s_{19}(g)=\kappa_3(19)(g)w\,s_O(g)
       \quad\text{in }N_2\rtimes G_{\mathbb Q}
 \tag{E.37.2}
\]
has no lift to $N_3\rtimes G_{\mathbb Q}$.
Here $\kappa_3(19)w$ is written additively
in the central subgroup
$\gamma_2N/\gamma_3N=\bigwedge^2V$.
The obstruction is detected by the actual
central coefficient pushout
$A\to\mathbb Z/9(2)$ induced by $\Psi$,
at the prime $19$, where its residue is
$6\in\mathbb Z/9$. Consequently no lift
to the geometrically pro-three metabelian
group, or to the full geometric group,
can have (E.37.2) as its class-two image.
\end{proposition}
\begin{proof}
Let $p:V\to T_3E$ be the actual compact
projection of E.18. Fix its determinant
pairing so that $p(w)=9$ in
$\bigwedge^2T_3E=\mathbb Z_3(1)$.
Use the global unit
\[
                 f=(X+Z)/(X+Y),
       \qquad\operatorname{div}(f)=3(B_0-O).
\]
Its geometric Kummer map
$\ell_f:V\to\mathbb Z_3(1)$ is integral
and vanishes on the rational compact
summand. Embed the free degree-three
Lie lattice in its tensor algebra and
apply the literal map
\[
 \Phi=(\ell_f\otimes\det(p,p)):
    A\lhook\joinrel\longrightarrow V^{\otimes3}
       \longrightarrow\mathbb Z_3(2).
 \tag{E.37.3}
\]
It is integral and Galois-equivariant.
The first-letter terms of $[z,[u,v]]$
give
\[
                  \Phi([z,w])=18\ell_f(z).
 \tag{E.37.4}
\]
The other two terms have first letter in
$H$, on which $\ell_f$ is zero. The same
factor count shows that $\Phi$ kills the
$HW^2$ multidegree.

Here is its complete integral row, so
that division by three is checked on
the actual lattice. In E.18's ordered
basis write $e_0=f_1$, $e_1=f_2$,
followed by the eight nonorigin cusps;
then $e_5=e_{B_0}$. The relation
$j_{0*}=\binom{3I}{U}$ gives
\[
 \ell_f=(0,2,0,0,0,3,0,0,0,0).
\]
On the Hall basis $[[e_i,e_j],e_k]$,
with $i>j$ and $k\ge j$, the only nonzero
values of $\Phi$ are
\[
 \begin{array}{c|rrr}
 (i,j,k)&(1,0,1)&(1,0,5)&(5,0,1)\\\hline
 \Phi&6&6&3\\
 \Psi&2&2&1
 \end{array}
 \tag{E.37.5}
\]
Thus $\Phi/3$ is an integral surjection;
after twisting by $(-1)$ it is exactly
the map $\Psi$ in the statement. The
division preserves equivariance because
the target lattice is torsion-free.
The program
\texttt{zero\_slice\_inverse\_limit.py}
constructs this row from the tensor
formula and checks $\Psi T=49\Psi$ and
$\Psi J=\Psi$ for the actual integral
Frobenius and conjugation matrices.

By E.35--E.36 the rational action class
consists of its evaluated bracket term
and an $HW^2$ elliptic term. The latter
is killed by (E.37.3). Since
$\ell_f(v_6/3+v_2/2)=1$, equations
(E.36.1) and (E.37.4) give
$\Psi_*[\eta]=6\kappa_3(3)$ rationally.
Both classes are unramified outside
three. The integral group
$H^1(G_{\mathbb Q,\{3,\infty\}},\mathbb Z_3(1))$
is the torsion-free module
$\mathbb Z_3\kappa_3(3)$, as proved in
E.27. The same equality therefore holds
integrally. Its possible coboundary is
$(\chi(g)-1)b$ for a fixed $b\in\mathbb Z_3(1)$.
Both cocycles vanish at complex
conjugation $c$: this is the normalization
of E.19 and the chosen real Kummer roots.
At $c$ the coboundary is $-2b$, hence
$b=0$. This proves the actual cocycle
identity (E.37.1).
At the normalized arithmetic generator
$\tau$ of E.31 its value is six.
Modulo nine a Tate coboundary vanishing
at $c$ is again zero. Thus the coefficient
class modulo nine is nonzero as claimed.

For (E.37.2), the equivariant inclusion
$\mathbb Z_3(1)w\subset\bigwedge^2V$
makes $\kappa_3(19)w$ an actual central
class-two cocycle. E.19 identifies its
connecting obstruction with
$[\eta\smile\kappa_3(19)]$.
Apply the map (E.37.1) to the first
factor. At $19$ its scalar tame
correction would have to be
\[
       -\frac{6\kappa_3(3)(F_{19})}{19-1}
          =-\frac{\kappa_3(3)(F_{19})}{3}.
 \tag{E.37.6}
\]
The unramified Kummer value has valuation
zero: $3^6\equiv7\pmod{19}$, so three
is not a cube in the residue field.
More precisely choose the compatible
root-of-unity coordinate with
$\zeta_9\equiv4\pmod{19}$. This element
has exact order nine, and
\[
         3^2\equiv9\equiv4^4\pmod{19}.
\]
The Frobenius Kummer value is therefore
$4\pmod9$. The right side of the scalar
equation before division is
$6\cdot4=6\pmod9$, whereas multiplication
by $19-1$ is zero modulo nine.
This is a finite coefficient obstruction
to the central pushout. In particular
the correction in (E.37.6) is not
integral, and no original lift exists.
The program verifies all of these finite
field residues.

Finally $N''\subset\gamma_4N$ gives the
continuous quotient map
$N/N''\twoheadrightarrow N_3$,
compatible with the genuine arithmetic
action and the map to $N_2$.
The full geometric group maps to this
pro-three quotient as well. Each proposed
lift along either map would give the
excluded class-three lift. These are
the comparison maps used in the last
assertion. The anchor is \S51.16,
folios 733v--736v, scans 172--178,
and LM7-E.23(i)--(ii).
\end{proof}

\begin{proposition}[E.38 --- Proved: every finite coefficient lifts, but the continuous class-three lift fails]
\label{E-zero-slice-inverse-limit}
In the same intrinsic cubic groups there
is an explicitly specified continuous
central class-two cocycle
\[
       z(g)=\beta(g)w,
 \qquad
       \beta=\sum_{n\ge1}3^{n+6}\kappa_3(p_n),
 \tag{E.38.1}
\]
with the distinct rational primes $p_n$
specified below by finite Frobenius
conditions, whose class-three connecting
class satisfies
\[
 \begin{gathered}
 0\ne[\eta\smile\beta]
             \in H^2_{\rm cts}(G_{\mathbb Q},A),\\
 [\eta\smile\beta]\longmapsto0
             \quad\text{in }H^2(G_{\mathbb Q},A/3^mA)
                    \quad\text{for every }m\ge1.
 \end{gathered}
 \tag{E.38.2}
\]
Equivalently, $z$ lifts through each
actual central extension
$N_3/(3^mA)\to N_2$, and has no
continuous lift through $N_3\to N_2$.
No compatible system of those finite
coefficient lifts is asserted. Every
finite partial sum in (E.38.1) does
have an integral continuous class-three
lift.
\end{proposition}
\begin{proof}
We specify the primes and all uniform
integrality bounds. Let $T_V$ be E.18's
actual Frobenius matrix at seven, after
the normalization in E.31. Thus the
arithmetic free generators satisfy
\[
 \chi(\tau)=\chi(\tau')=7,
 \quad \kappa_3(3)(\tau)=1,
 \quad \kappa_3(3)(\tau')=-1,
 \quad \tau'=c\tau c.
 \tag{E.38.3}
\]
Their actions on $V$ commute: over
$K=\mathbb Q(\zeta_3)$ the compact CM
image is abelian and the boundary image
is cyclotomic, and the rational compact
splitting is equivariant. Equality of
the rational matrices is equality of
these integral matrices. Set
$r=[\tau,\tau']=\tau\tau'\tau^{-1}\tau'^{-1}$.
The cocycle law gives exactly
\[
           \rho_V(r)=1,\qquad
           \chi(r)=1,\qquad
           \kappa_3(3)(r)=-12.
 \tag{E.38.4}
\]
For $n\ge1$ define
\[
                  g_n=r\tau^{3^{n+1}}.
 \tag{E.38.5}
\]
These are elements of the actual
arithmetic quotient $\Gamma$.
Their finite images are thus in the
image of $G_{\mathbb Q}$.
The geometric-series cocycle formula
and the elementary lifting-the-exponent
calculation give
\[
 \begin{aligned}
 \rho_V(g_n)&=T_V^{3^{n+1}},\\
 \chi(g_n)&=7^{3^{n+1}},
       &v_3(\chi(g_n)-1)&=n+2,\\
 \kappa_3(3)(g_n)&=-12+
       \frac{7^{3^{n+1}}-1}{6},
       &v_3(\kappa_3(3)(g_n))&=1.
 \end{aligned}
 \tag{E.38.6}
\]
For the last equality the second term
has valuation $n+1\ge2$, while the
first has valuation one.

Here is the fixed matrix certificate
needed for every $n$. On $A(-1)$ put
\[
 R=7^{-1}\operatorname{Lie}_3(T_V),
       \qquad D=R^3-I.
\]
The actual residual $T_V$ is unipotent
of exponent three, so $D\in3M_{330}(\mathbb Z_3)$.
Its complete Smith valuation list is
\[
                 1^{132},2^{82},3^{52},
                       4^{49},5^{14},6^1.
 \tag{E.38.7}
\]
The executed program constructs every
entry, finds all $330$ pivots at
precision $3^{12}$, and independently
checks determinant valuation $724$.
For that second check put
$\pi=1+3\zeta$, so
$\pi+\bar\pi=-1$, $\pi\bar\pi=7$.
The contributions to the determinant
valuation from the five eigenvalue
blocks of $R^3-I$ are
\[
                  4+8\cdot4+16\cdot2
                       +64\cdot5+168\cdot2=724.
\]
These follow respectively from
$\operatorname{Norm}(\pi^3-1)$,
$\operatorname{Norm}(\pi^6-1)$,
$7^3-1$, $\operatorname{Norm}((7\pi)^3-1)$,
and $7^6-1$; the multiplicities are
the $H^3,H^2W,HW^2,W^3$ multidegrees.
All pivots lie below the working
precision, so the reported valuations
lift under integral row and column
operations. In particular $D$ is
invertible rationally and $3^6D^{-1}$
is integral.

Define
$D_n=(I+D)^{3^n}-I=R^{3^{n+1}}-I$.
For any matrix $B\in3M_{330}(\mathbb Z_3)$,
\[
 (I+B)^3-I=3B(I+B+B^2/3),
          \qquad I+B+B^2/3\in I+3M_{330}(\mathbb Z_3).
\]
Induction gives $D_n=3^nD U_n$ for an
integral invertible matrix $U_n$.
Therefore
\[
                        3^{n+6}D_n^{-1}
                             \in M_{330}(\mathbb Z_3).
 \tag{E.38.8}
\]
This proves the bound at every depth,
not merely at the computed depth.

To choose $p_n$, take the finite image
of the actual homomorphism consisting
of $\rho_V$, $\chi$, and the affine
Kummer representation $(\kappa_3(3),\chi)$,
all reduced modulo $3^{n+8}$.
The last pair is a homomorphism into
the corresponding affine semidirect
product by the cocycle law. Let $p_n$
be the least prime larger than $p_{n-1}$
and different from three whose Frobenius
conjugacy class in this finite image
is that of $g_n$; set $p_0=3$.
Chebotarev gives infinitely many such
primes at each step, so this specifies
the entire sequence. Choose a
prolongation with representative equal
to the chosen finite image of $g_n$.
The condition modulo three puts $p_n$
in the split class of $K/\mathbb Q$.
The two primes of $K$ above $p_n$ are
related by integral complex conjugation;
the estimates below hold at both.
The specification uses only the named
finite arithmetic representations and
their actual Frobenius classes. We have
not enumerated these primes numerically.

Write
$D_{p_n}=p_n^{-1}\operatorname{Lie}_3
(\rho_V(F_{p_n}))-I$.
The finite Frobenius condition gives
$D_{p_n}-D_n\in3^{n+8}M_{330}(\mathbb Z_3)$.
By (E.38.8),
$D_n^{-1}(D_{p_n}-D_n)\in9M_{330}(\mathbb Z_3)$.
The Neumann series makes the resulting
$I+9M$ matrix invertible integrally.
Thus the inverse bound (E.38.8) also
holds with $D_{p_n}$, and (E.38.6)
gives the exact valuations
\[
 v_3(p_n-1)=n+2,
 \qquad v_3(\kappa_3(3)(F_{p_n}))=1.
 \tag{E.38.9}
\]

The sum (E.38.1) converges uniformly
as a function on $G_{\mathbb Q}$:
every normalized integral Kummer
cocycle has values in $\mathbb Z_3$,
and $3^{n+6}$ tends to zero. Its limit
is therefore a continuous Tate cocycle,
and $z=\beta w$ is an actual central
class-two cocycle. The unique rational
tame correction prescribed by E.19
at $p_n$ is
\[
 x_n=-3^{n+6}D_{p_n}^{-1}\eta(F_{p_n}).
 \tag{E.38.10}
\]
It is integral by the proved matrix
bound. All other primes outside three
have zero residue and zero correction.
Nevertheless these integral corrections
do not tend to zero. Equivariance of
E.37's integral map says
$\Psi D_{p_n}=(p_n-1)\Psi$; hence
\[
 \Psi x_n=-\frac{3^{n+6}\,6\kappa_3(3)(F_{p_n})}
                            {p_n-1},
       \qquad v_3(\Psi x_n)=6
                  \quad\text{for every }n.
 \tag{E.38.11}
\]
An integral linear map takes a sequence
tending to zero to one tending to
zero. Formula (E.38.11) therefore
violates precisely the restricted
continuity condition of (E19.5).
That exact criterion proves the first
nonvanishing in (E.38.2).

For a finite partial sum of $\beta$,
(E.38.10) is integral and has finite
support, so it satisfies the same
criterion and gives an integral
continuous class-three lift. For
fixed $m$, all terms with $n+6\ge m$
vanish modulo $3^m$. The connecting
cocycle $\eta\smile\beta$ modulo
$3^mA$ therefore equals the reduction
of the obstruction of one finite
partial sum. The latter is zero.
Central-extension obstruction theory
then gives the stated lift through
$N_3/(3^mA)$ for each $m$.
The preceding nonzero continuous
obstruction proves that these
existence statements do not supply
a compatible inverse-limit lift.
This is the intrinsic cubic
zero-abelian instance of the exact
finite-coefficient comparison asked
in LM7-E.23(i), with its actual
arithmetic action retained.
\end{proof}
